% Options for packages loaded elsewhere
\PassOptionsToPackage{unicode}{hyperref}
\PassOptionsToPackage{hyphens}{url}
\documentclass[
  ignorenonframetext,
]{beamer}
\newif\ifbibliography
\usepackage{pgfpages}
\setbeamertemplate{caption}[numbered]
\setbeamertemplate{caption label separator}{: }
\setbeamercolor{caption name}{fg=normal text.fg}
\beamertemplatenavigationsymbolsempty
% remove section numbering
\setbeamertemplate{part page}{
  \centering
  \begin{beamercolorbox}[sep=16pt,center]{part title}
    \usebeamerfont{part title}\insertpart\par
  \end{beamercolorbox}
}
\setbeamertemplate{section page}{
  \centering
  \begin{beamercolorbox}[sep=12pt,center]{section title}
    \usebeamerfont{section title}\insertsection\par
  \end{beamercolorbox}
}
\setbeamertemplate{subsection page}{
  \centering
  \begin{beamercolorbox}[sep=8pt,center]{subsection title}
    \usebeamerfont{subsection title}\insertsubsection\par
  \end{beamercolorbox}
}
% Prevent slide breaks in the middle of a paragraph
\widowpenalties 1 10000
\raggedbottom
\AtBeginPart{
  \frame{\partpage}
}
\AtBeginSection{
  \ifbibliography
  \else
    \frame{\sectionpage}
  \fi
}
\AtBeginSubsection{
  \frame{\subsectionpage}
}
\usepackage{iftex}
\ifPDFTeX
  \usepackage[T1]{fontenc}
  \usepackage[utf8]{inputenc}
  \usepackage{textcomp} % provide euro and other symbols
\else % if luatex or xetex
  \usepackage{unicode-math} % this also loads fontspec
  \defaultfontfeatures{Scale=MatchLowercase}
  \defaultfontfeatures[\rmfamily]{Ligatures=TeX,Scale=1}
\fi
\usepackage{lmodern}
\ifPDFTeX\else
  % xetex/luatex font selection
\fi
% Use upquote if available, for straight quotes in verbatim environments
\IfFileExists{upquote.sty}{\usepackage{upquote}}{}
\IfFileExists{microtype.sty}{% use microtype if available
  \usepackage[]{microtype}
  \UseMicrotypeSet[protrusion]{basicmath} % disable protrusion for tt fonts
}{}
\makeatletter
\@ifundefined{KOMAClassName}{% if non-KOMA class
  \IfFileExists{parskip.sty}{%
    \usepackage{parskip}
  }{% else
    \setlength{\parindent}{0pt}
    \setlength{\parskip}{6pt plus 2pt minus 1pt}}
}{% if KOMA class
  \KOMAoptions{parskip=half}}
\makeatother
\setlength{\emergencystretch}{3em} % prevent overfull lines
\providecommand{\tightlist}{%
  \setlength{\itemsep}{0pt}\setlength{\parskip}{0pt}}
\usepackage{bookmark}
\IfFileExists{xurl.sty}{\usepackage{xurl}}{} % add URL line breaks if available
\urlstyle{same}
\hypersetup{
  hidelinks,
  pdfcreator={LaTeX via pandoc}}

\author{\texorpdfstring{}{}}
\date{}

\begin{document}

\begin{frame}[fragile]{Introduction and Scope}
\protect\phantomsection\label{introduction-and-scope}
Machine learning on source code has matured alongside static analysis:
both need structured representations of programs, but learning pipelines
additionally need fixed feature layouts, batchable graphs, and
reproducible exports {[}@allamanis2018survey{]}. COGANT addresses that
gap by making the path from repository checkout to \textbf{Generalized
Notation Notation (GNN) bundles} explicit, inspectable, and
configurable.

\begin{block}{Terminology: ``GNN'' in this manuscript}
\protect\phantomsection\label{terminology-gnn-in-this-manuscript}
In this manuscript, \textbf{GNN} refers to the Active Inference
Institute's \textbf{Generalized Notation Notation}
(\href{https://github.com/ActiveInferenceInstitute/GeneralizedNotationNotation}{repository})
{[}@friedman2024gnn{]}, a structured notation for state-space and
process models---\emph{not} graph neural networks. COGANT translates
source code into this notation, producing program-graph and state-space
artifacts (state space, observation modalities, actions/policies,
likelihood structure, preferences, time settings, parameterization) that
downstream tools can consume. Those downstream tools include graph
neural network training pipelines, which can ingest the exported tensor
views of the program graph, but such neural networks are a
\emph{consumer} of COGANT's output, not its output format. Where the
manuscript discusses graph neural networks as related work, it uses the
phrase ``graph neural networks'' in full; the acronym GNN always refers
to Generalized Notation Notation.
\end{block}

\begin{block}{Problem}
\protect\phantomsection\label{problem}
A research or engineering group typically has:

\begin{itemize}
\tightlist
\item
  \textbf{Source code} in one or more languages.
\item
  \textbf{Analysis goals} that may mix classical queries (who calls
  whom?) with learned models (bug detection, retrieval, summarization).
\item
  \textbf{Tooling fragmentation}: compilers and linters produce one
  shape of facts; PyTorch or JAX expect another.
\end{itemize}

COGANT unifies these behind a single pipeline whose intermediate
artifacts are documented as a progression of IRs (repo IR, program
graph, semantic mapping, state space, process model, validation), as
described in \texttt{../cogant/docs/ARCHITECTURE.md} and
\texttt{../cogant/docs/SPEC.md}.
\end{block}

\begin{block}{What COGANT is not}
\protect\phantomsection\label{what-cogant-is-not}
COGANT is not a replacement for a full compiler IR, a theorem prover, or
a security scanner by itself. It \textbf{extracts and serializes}
program graphs and behavioral sketches for downstream use. Language
coverage and rule depth follow the implementation-status table in the
SPEC: Python is the primary front end at v0.1.x; additional languages
and Rust acceleration are staged as documented there.
\end{block}

\begin{block}{Manuscript versus package docs}
\protect\phantomsection\label{manuscript-versus-package-docs}
The canonical API and CLI details live in:

\begin{itemize}
\tightlist
\item
  \texttt{../cogant/docs/API\_GUIDE.md} --- Session,
  \texttt{PipelineRunner}, \texttt{Bundle}, \texttt{ReviewAPI}
\item
  \texttt{../cogant/docs/CLI\_GUIDE.md} --- commands and flags
\item
  \texttt{../cogant/docs/GNN\_EXPORT.md} --- Generalized Notation
  Notation schema, section ordering, and companion interop field
  contracts (including optional PyG/DGL tensor views)
\item
  \texttt{../cogant/docs/PLUGIN\_API.md} --- parsers, rules, validators,
  exporters
\item
  \texttt{../cogant/docs/ARCHITECTURE.md} --- layers, crates, data flow
\end{itemize}

This manuscript summarizes theory and practice for readers who want a
single linear narrative before diving into those files.
\end{block}

\begin{block}{Dual Python/Rust architecture}
\protect\phantomsection\label{dual-pythonrust-architecture}
COGANT employs a dual-language architecture. The \textbf{Python
orchestration layer} handles session management, pipeline coordination,
configuration, file discovery, AST parsing, and the plugin interface --
tasks where developer ergonomics and extensibility matter more than raw
throughput. The \textbf{Rust core}, organized as a workspace of seven
crates (\texttt{cogant-core}, \texttt{cogant-graph},
\texttt{cogant-translate}, \texttt{cogant-statespace},
\texttt{cogant-gnn}, \texttt{cogant-store}, \texttt{cogant-ffi}),
implements typed graph operations, translation engine internals,
state-space structures, and Generalized Notation Notation (GNN)
formatting -- tasks where memory layout and cache locality dominate
wall-clock time. Communication between layers flows through PyO3
bindings (\texttt{cogant-ffi}), which release the Python GIL during Rust
execution to permit concurrent stage processing. Where Rust bindings are
not yet wired for a code path (see the implementation-status table in
\texttt{../cogant/docs/SPEC.md}), Python fallback implementations ensure
the pipeline remains functional at the cost of higher latency on large
repositories.
\end{block}

\begin{block}{Roadmap of the following sections}
\protect\phantomsection\label{roadmap-of-the-following-sections}
Section 2 formalizes the program graph and IR progression. Section 3
describes the realized pipeline behavior and artifacts. Sections 4--7
cover conclusions, how to run experiments, reproducibility, and related
work.
\end{block}
\end{frame}

\end{document}
